\section{Introduction and main results}

A periodic signing of a cycle square can be studied through a single cell
and its boundary phase. This decomposition depends on repeating the same
cell. We ask what spectral control remains when the cell lengths are
chosen independently within one cyclic concatenation.

For $N\ge8$, the cycle square $C_N(1,2)$ has vertex set $\Z/N\Z$ and
edges $\{i,i+1\}$ and $\{i,i+2\}$. A real signing assigns $\pm1$ to each
edge. For its symmetric adjacency matrix $A$, write
\[
 \rho(A)=\max\{|\lambda|:\lambda\in\spec(A)\},
 \qquad \mu_N=\min_\sigma\rho(A_\sigma).
\]
The minimization controls both spectral edges. The benchmark in the
original twisted-optimality question \cite{Suvagiya2026old} is
\begin{equation}
 \tw(N)^2=4+2\cos\frac{2\pi}{N}+2\cos\frac{4\pi}{N}
 \qquad(N\text{ even},\ N\ge8).
 \label{eq:twisted}
\end{equation}

Our cells have a fixed internal pattern and fixed ends. They are built
from the period-eight word
\begin{equation}
 t=(1,1,-1,1,-1,-1,1,-1),\qquad
 w_j=t^j\mathbin{\|}(1,-1),\qquad h_j=8j+2\quad(j\ge1).
 \label{eq:word}
\end{equation}
Here $\|$ denotes concatenation. We join any positive number of these
triangle-sign words, choosing the integers $j\ge1$ independently.
Appending $(1,-1)$ adds two vertices to the period-eight bulk, so that
$r$ such cells give order $2r$ modulo eight. In the associated
quadrilateral signs $Q_u=\tau_u\tau_{u+1}$, the positive positions are
separated by four within a cell and by six across a boundary. The precise
words, rather than the gap count alone, are the hypotheses.

Two uniformities are required. The lengths must be allowed to vary
independently within a single concatenation, and the estimates must remain
controlled when arbitrarily many cells are joined. Summing one error per
cell would lose the second uniformity. The equal-cell Fourier
decomposition does not directly address the first.

Both difficulties can be resolved at the retained junctions. After the
interiors are eliminated, the diagonal block at each junction combines the
responses of its two adjacent cells. At a fixed spectral cap, these blocks
approach the same positive six-dimensional matrix as the adjacent cells
become long. The remaining interactions join neighboring junctions, and
their quadratic-form estimate counts two chain ends at each junction.
Thus the shortest cell controls the local error, while the cyclic
incidence keeps its assembly independent of the number of cells.

For a length-$N$ triangle-sign word
$\tau$ and $\alpha\in\{\pm1\}$, define $A_N(\tau,\alpha)$ by the edge signs
\begin{equation}
 a_i=\begin{cases}1&i<N-1,\\\alpha&i=N-1,\end{cases}
 \quad
 \sigma\{i,i+1\}=a_i,\qquad
 \sigma\{i,i+2\}=\tau_i a_i a_{i+1},
 \label{eq:gauge-definition}
\end{equation}
with cyclic indices. The triangle signs are $\tau_i$, and the
Hamilton-cycle sign product is $\alpha$.
The two rational caps below are certified upper thresholds for these
constructions, rather than proposed values of the unrestricted minimum:
define
\[
 \czero=\frac{198}{25}=7.92,\qquad
 \cone=\frac{790537}{100000}=7.90537.
\]

\begin{theorem}[Unequal-cell bound]\label{thm:cells}
Let $r\ge1$, $j_0,\ldots,j_{r-1}\ge1$, and
$\tau=w_{j_0}\|\cdots\|w_{j_{r-1}}$, of total length $N$.
For either holonomy $\alpha$,
\begin{align}
 \min_i h_{j_i}\ge106&\quad\Longrightarrow\quad
       \rho(A_N(\tau,\alpha))^2<\czero,\label{eq:cell-cap0}\\
 \min_i h_{j_i}\ge202&\quad\Longrightarrow\quad
       \rho(A_N(\tau,\alpha))^2<\cone.\label{eq:cell-cap1}
\end{align}
The constants do not depend on $r$, and the lengths need not be equal.
\end{theorem}

The bounds hold for both holonomies and for independently chosen cell
lengths. They control both spectral edges through $\rho(A)^2$.
The length thresholds are sufficient, and neither they nor the spectral
caps are asserted to be optimal. The prescribed words remain essential:
the theorem does not cover arbitrary signings or arbitrary arrangements
of local defects.

The single-cell family permits a sharper finite completion.
\begin{theorem}[One-cell bound]\label{thm:one-cell}
For every $k\ge1$,
\[
 \rho(A_{8k+2}(w_k,+1))^2<\cone.
\]
Moreover,
\begin{equation}
 \frac{7905369}{1000000}
 <\sup_{k\ge1}\rho(A_{8k+2}(w_k,+1))^2
 \le\frac{790537}{100000}.
 \label{eq:supremum}
\end{equation}
\end{theorem}
The interval in \eqref{eq:supremum} has width $10^{-6}$. Its upper endpoint
is non-strict; no convergence, monotonicity, or exact supremum is asserted.

To cover each nonzero even residue modulo eight, put
\[
 W_{k,r}=\mathop{\|}_{a=0}^{r-1}w_{\lfloor(k+a)/r\rfloor},
 \qquad r\in\{1,2,3\},\ k\ge6.
\]
The cell parameters sum to $k$, so the word has length $8k+2r$.
\begin{theorem}[Explicit finite competitors]\label{thm:witnesses}
For $r\in\{1,2,3\}$ and $k\ge6$, the matrix
\[
 B_{k,r}=A_{8k+2r}(W_{k,r},(-1)^{r+1})
\]
satisfies
\begin{equation}
 \rho(B_{k,r})^2<\czero<\tw(8k+2r)^2.
 \label{eq:nonzero-witnesses}
\end{equation}
For $r=1$ the stronger cap $\cone$ holds. Together with the period-eight
family, these constructions prove
\[
 \mu_N<\tw(N)
 \quad\text{for }N\in\{32,40\}\cup\{N\ge48:N\text{ even}\}.
\]
\end{theorem}

The period-eight predecessor \cite{EarlierPeriodEight} already supplies the
orders $32$ and $40$, and the earlier Target A package
\cite{EarlierTargetA} supplies every even order at least $48$.
These are unpublished project predecessors. The contribution here is a
common analytic mechanism, sharper quantitative bounds, and a finite
completion with directly specified signings. For odd $k$ in residue four, and for
$3\nmid k$ in residue six, our placements need not be the older balanced-gap
placements. Their finite certificates are recomputed for the words above.
We do not reprove the universal lower bounds at complementary smaller
orders, determine $\mu_N$ when the twisted signing fails, or classify
minimizing signings.

The literature target has also changed. The original preprint
\cite{Suvagiya2026old} was withdrawn and merged into
\cite{Suvagiya2026}. Its revised Theorem 26 supplies the positive-holonomy
period-eight value
\begin{equation}
 \etaedge=4+\sqrt{10+2\sqrt5},\qquad \rstar=\sqrt{\etaedge},
 \label{eq:rstar}
\end{equation}
for $N=8m$, $m\ge4$, while Conjecture 28 proposes $\mu_{8m}=\rstar$
throughout that range. The exact negative-holonomy formula from the earlier
project manuscript \cite{EarlierPeriodEight} is strictly smaller at every
finite size. We provide its full proof and apply it to that revised claim.
This application is distinct from a priority claim for the inherited formula.

Schur complements and contractive Riccati recurrences are standard matrix
tools. The contribution here is their quantitative assembly for the stated
cells: a common limiting core controls unequal chains, and the retained
matrix has two chain incidences per boundary. Related signed-spectral
results concern different questions. Belardo and Brunetti
\cite[Theorem 1.3 and Lemma 3.8]{BelardoBrunetti2024} determine limit points
of signed adjacency spectral radii and use a Schur-complement calculation
for an auxiliary family. Our one-cell enclosure does not identify a new
limit point or prove convergence. Kannan and Pragada
\cite[Theorem 3.3]{KannanPragada2023} bound the largest adjacency eigenvalue
using edge count, frustration index, and balanced clique number. Here the
input is the prescribed triangle-sign word and its boundary response; the
result controls both spectral edges uniformly in cell length and count.

The use of holonomy belongs to the established character approach to
spectral graph theory; Luo and Roy \cite[Theorems 3.6 and 3.13]{LuoRoy2026}
study character families and endpoint attainment. Those results do not
supply the fixed-support minimum here. Likewise, signed moments averaged
over all signings, as in Chen, van Dam, and Bu
\cite[Theorem 3.1]{ChenDamBu2024}, are not pointwise lower bounds for every
signing. Our estimates instead arise from an explicit fiber calculation
and exact positive-definiteness tests after a uniform analytic reduction.
The latter proofs are finite-certificate-assisted, with independent exact
replays; their finite checks are specified in Section~\ref{sec:finite} and
located in the frozen supplement \cite{C029Supplement}.

\paragraph{Proof overview.}
We prove Theorem~\ref{thm:cells} by testing the positivity of $cI-A^2$.
The prescribed cell ends allow us to retain six coordinates at every
junction and eliminate the intervening four-site blocks. Their Schur
complements follow one length-independent Riccati recurrence. A
contraction estimate controls the pivots, while separate response
estimates control the contributions of the eliminated interior to its
two ends. Comparing a finite seed with its bounded tail proves positivity
of the limiting junction matrix. This argument is carried out separately
at each of the two spectral caps.

For unequal cells, the finite responses are assembled before any limit
is taken. Relative to a block diagonal sum of the limiting matrices, the
diagonal errors are controlled by the shortest cell. The cross terms
satisfy a quadratic-form bound in which each junction is counted twice,
including the one-cell loop and the two-cell parallel contributions.
Its constant is therefore independent of the number of cells.
Positivity of the assembled boundary matrix, together with positivity of
the eliminated pivots, gives $cI-A^2\succ0$ by Schur congruence.
Exact rational checks establish the finite inequalities needed to enter
the contractive regime; the analytic estimates extend them to all
admissible lengths and cell counts.

Sections~\ref{sec:chain}--\ref{sec:cells} carry out this argument, and
Figure~\ref{fig:boundary} displays the boundary reduction.
Section~\ref{sec:finite} completes the bounded ranges for the one-cell
bound and the two- and three-cell competitors. The residue-zero
competitors come from the independent period-eight calculation in
Sections~\ref{sec:holonomy}--\ref{sec:period8}, which also gives the
application to the revised conjecture. Thus the full witness range uses
both proof tracks.
